\documentclass[10pt]{article}
\usepackage{tp2style}
\hypersetup{pdftitle={Requirements to Tests Map},pdfauthor={CSE 360 Team 24}}
\fancyfoot[L]{\sffamily\scriptsize\color{muted}Requirements to Tests Map \textperiodcentered{} Team 24}
\newcommand{\nauto}{76}\newcommand{\nman}{7}\newcommand{\ntotal}{83}

\begin{document}
\doctitle{Requirements to Tests Map}{\team \\ First issued October 22, 2026; this revision October 26, 2026}

\section*{What this map covers}
CS-1 to CS-12 are the Contributor user stories split into requirements that can each be tested
on their own (agreed by the team on October 13). Every requirement has at least one positive
test, showing it is met for valid use, and at least one negative test, showing incorrect use is
refused or explained. There are \ntotal{} tests: \nauto{} semi-automated tests in four suites,
run together by \texttt{TP2TestRunner}, and \nman{} manual procedures (MAN) for what only a
window can show. Every test is one method; its Javadoc states the requirement, how the test
works, and what output means a pass. The HTML pages and screenshots named here are in the
\texttt{ContributorTests\_Javadocs} folder.

\subsection*{How to run the tests}
\begin{tabularx}{\linewidth}{@{}l X@{}}
Automated & Run \texttt{testingAutomation.TP2TestRunner} (Eclipse: Run As, Java Application, with
  the JavaFX module path and the H2 jar as for \texttt{FoundationsMain}). Each test prints one line
  starting \texttt{Pass:} or \texttt{Failure:} followed by its id and requirement; the last line is
  \texttt{RESULT: every test behaved as expected} only when nothing failed. \\
Manual & Run \texttt{testingAutomation.LessonsManualTestProcedures} beside a running
  \texttt{FoundationsMain}. It prints each procedure's steps and expected result and records
  P, F, or S for each. \\
\end{tabularx}

\section*{Requirements and their tests}
\begin{small}
\setlength{\tabcolsep}{4pt}\begin{longtable}{@{}>{\raggedright\arraybackslash}p{0.055\linewidth} >{\raggedright\arraybackslash}p{0.33\linewidth} >{\raggedright\arraybackslash}p{0.14\linewidth} >{\raggedright\arraybackslash}p{0.15\linewidth} >{\raggedright\arraybackslash}p{0.245\linewidth}@{}}
\toprule
\textbf{Req.} & \textbf{Requirement} & \textbf{Positive tests} & \textbf{Negative tests} & \textbf{HTML pages} \\
\midrule
\endhead
\textbf{CS-1} & A contributor can create a lesson with a title, category, description, and optional resolution and keywords. The contributor is set from the logged-in user. & CRUD-01, MOD-01, MAN-02 & CRUD-02, MOD-02 & \scriptsize LessonsCrudTestingAutomation.html \newline LessonsManualTestProcedures.html \newline LessonsModelTestingAutomation.html \\ \addlinespace[3pt]
\textbf{CS-2} & A contributor can see all lessons, newest first. An empty list is explained in words, not shown as a blank screen. & LIST-01, LIST-11, MOD-03 & LIST-02, MOD-19, MAN-01 & \scriptsize LessonsListTestingAutomation.html \newline LessonsManualTestProcedures.html \newline LessonsModelTestingAutomation.html \\ \addlinespace[3pt]
\textbf{CS-3} & A contributor can open one lesson and see all of its fields. & CRUD-03 & CRUD-04 & \scriptsize LessonsCrudTestingAutomation.html \\ \addlinespace[3pt]
\textbf{CS-4} & A contributor can search lessons. The search covers title, description, and keywords and ignores case. No match is explained in words. & LIST-03, LIST-10, LIST-12, MOD-05 & LIST-04, LIST-13, MOD-06, MOD-17 & \scriptsize LessonsListTestingAutomation.html \newline LessonsModelTestingAutomation.html \\ \addlinespace[3pt]
\textbf{CS-5} & A contributor can filter lessons by category. An empty category is explained in words. The filter's label “All categories” cannot be used as a category. & LIST-05, LIST-09, LIST-10, MOD-07, MOD-24 & LIST-06, MOD-08, MOD-18, MOD-23 & \scriptsize LessonsListTestingAutomation.html \newline LessonsModelTestingAutomation.html \\ \addlinespace[3pt]
\textbf{CS-6} & A contributor can view only their own lessons. If they have none, a message says so. & LIST-07, MOD-09 & LIST-08, MOD-04, MOD-10 & \scriptsize LessonsListTestingAutomation.html \newline LessonsModelTestingAutomation.html \\ \addlinespace[3pt]
\textbf{CS-7} & A contributor can update a lesson they contributed. An update is validated the same way as a new lesson. & CRUD-05, MOD-13 & CRUD-06, MOD-14 & \scriptsize LessonsCrudTestingAutomation.html \newline LessonsModelTestingAutomation.html \\ \addlinespace[3pt]
\textbf{CS-8} & A contributor can delete a lesson they contributed, after confirming. It is removed for all users. & CRUD-07, MOD-15, MAN-04, MAN-06 & CRUD-08, MOD-16, MAN-04 & \scriptsize LessonsCrudTestingAutomation.html \newline LessonsManualTestProcedures.html \newline LessonsModelTestingAutomation.html \\ \addlinespace[3pt]
\textbf{CS-9} & Only the contributor of a lesson can update or delete it. Any other user is refused with a message. & CRUD-09, MOD-11 & CRUD-10, MOD-12, MOD-16, MOD-20, MAN-05 & \scriptsize LessonsCrudTestingAutomation.html \newline LessonsManualTestProcedures.html \newline LessonsModelTestingAutomation.html \\ \addlinespace[3pt]
\textbf{CS-10} & Title, category, description, and contributor are required; resolution and keywords are optional. A missing required value is refused with a message naming the field. & VAL-01, VAL-06, VAL-10, VAL-15, VAL-18, VAL-22, VAL-24 & VAL-02, VAL-03, VAL-07, VAL-11, VAL-23, VAL-25, MAN-03 & \scriptsize LessonValidationTestingAutomation.html \newline LessonsManualTestProcedures.html \\ \addlinespace[3pt]
\textbf{CS-11} & Length limits: title 80, category 30, description 2000, resolution 2000, keywords 200 characters, at most 10 keywords of 30 characters each. A value over a limit, or an empty keyword entry, is refused with a message naming the problem. & VAL-04, VAL-08, VAL-12, VAL-13, VAL-16, VAL-19, VAL-24 & VAL-05, VAL-09, VAL-14, VAL-17, VAL-20, VAL-21, MOD-21, MOD-22 & \scriptsize LessonValidationTestingAutomation.html \newline LessonsModelTestingAutomation.html \\ \addlinespace[3pt]
\textbf{CS-12} & Lessons persist across application restarts, including on a database left by an HW2 version. & CRUD-11, CRUD-13, MAN-07 & CRUD-12, CRUD-14 & \scriptsize LessonsCrudTestingAutomation.html \newline LessonsManualTestProcedures.html \\ \addlinespace[3pt]
\bottomrule\end{longtable}
\end{small}

\section*{Every test, its Javadoc, and its screenshot}
The anchor is the method name inside the HTML page: \texttt{\#crud01CreateValidLesson()} in
\texttt{LessonsCrudTestingAutomation.html} opens that test's Javadoc. CRUD-11 and CRUD-12, and CRUD-13 and CRUD-14, share one method because
the negative test checks the database the positive test reopened.

\begin{small}
\setlength{\tabcolsep}{4pt}\begin{longtable}{@{}>{\raggedright\arraybackslash}p{0.07\linewidth} >{\raggedright\arraybackslash}p{0.08\linewidth} >{\raggedright\arraybackslash}p{0.07\linewidth} >{\raggedright\arraybackslash}p{0.33\linewidth} >{\raggedright\arraybackslash}p{0.37\linewidth}@{}}
\toprule
\textbf{Test} & \textbf{Req.} & \textbf{Kind} & \textbf{Method (anchor in the HTML page)} & \textbf{Screenshot} \\
\midrule
\endhead
\multicolumn{5}{@{}l}{\sffamily\bfseries\small LessonValidationTestingAutomation.html} \\
VAL-01 & CS-10 & pos. & \scriptsize\texttt{\#val01TitleAccepted()} & \scriptsize\texttt{7.4-VAL-01-TitleAccepted.png} \\
VAL-02 & CS-10 & neg. & \scriptsize\texttt{\#val02BlankTitleRefused()} & \scriptsize\texttt{7.4-VAL-02-BlankTitleRefused.png} \\
VAL-03 & CS-10 & neg. & \scriptsize\texttt{\#val03WhitespaceTitleRefused()} & \scriptsize\texttt{7.4-VAL-03-WhitespaceTitleRefused.png} \\
VAL-04 & CS-11 & pos. & \scriptsize\texttt{\#val04TitleAt80Accepted()} & \scriptsize\texttt{7.4-VAL-04-TitleAt80Accepted.png} \\
VAL-05 & CS-11 & neg. & \scriptsize\texttt{\#val05TitleAt81Refused()} & \scriptsize\texttt{7.4-VAL-05-TitleAt81Refused.png} \\
VAL-06 & CS-10 & pos. & \scriptsize\texttt{\#val06CategoryAccepted()} & \scriptsize\texttt{7.4-VAL-06-CategoryAccepted.png} \\
VAL-07 & CS-10 & neg. & \scriptsize\texttt{\#val07BlankCategoryRefused()} & \scriptsize\texttt{7.4-VAL-07-BlankCategoryRefused.png} \\
VAL-08 & CS-11 & pos. & \scriptsize\texttt{\#val08CategoryAt30Accepted()} & \scriptsize\texttt{7.4-VAL-08-CategoryAt30Accepted.png} \\
VAL-09 & CS-11 & neg. & \scriptsize\texttt{\#val09CategoryAt31Refused()} & \scriptsize\texttt{7.4-VAL-09-CategoryAt31Refused.png} \\
VAL-10 & CS-10 & pos. & \scriptsize\texttt{\#val10DescriptionAccepted()} & \scriptsize\texttt{7.4-VAL-10-DescriptionAccepted.png} \\
VAL-11 & CS-10 & neg. & \scriptsize\texttt{\#val11BlankDescriptionRefused()} & \scriptsize\texttt{7.4-VAL-11-BlankDescriptionRefused.png} \\
VAL-12 & CS-11 & pos. & \scriptsize\texttt{\#val12DescriptionAt300Accepted()} & \scriptsize\texttt{7.4-VAL-12-DescriptionAt300Accepted.png} \\
VAL-13 & CS-11 & pos. & \scriptsize\texttt{\#val13DescriptionAt2000Accepted()} & \scriptsize\texttt{7.4-VAL-13-DescriptionAt2000Accepted.png} \\
VAL-14 & CS-11 & neg. & \scriptsize\texttt{\#val14DescriptionAt2001Refused()} & \scriptsize\texttt{7.4-VAL-14-DescriptionAt2001Refused.png} \\
VAL-15 & CS-10 & pos. & \scriptsize\texttt{\#val15EmptyResolutionAccepted()} & \scriptsize\texttt{7.4-VAL-15-EmptyResolutionAccepted.png} \\
VAL-16 & CS-11 & pos. & \scriptsize\texttt{\#val16ResolutionAt2000Accepted()} & \scriptsize\texttt{7.4-VAL-16-ResolutionAt2000Accepted.png} \\
VAL-17 & CS-11 & neg. & \scriptsize\texttt{\#val17ResolutionAt2001Refused()} & \scriptsize\texttt{7.4-VAL-17-ResolutionAt2001Refused.png} \\
VAL-18 & CS-10 & pos. & \scriptsize\texttt{\#val18EmptyKeywordsAccepted()} & \scriptsize\texttt{7.4-VAL-18-EmptyKeywordsAccepted.png} \\
VAL-19 & CS-11 & pos. & \scriptsize\texttt{\#val19ValidKeywordsAccepted()} & \scriptsize\texttt{7.4-VAL-19-ValidKeywordsAccepted.png} \\
VAL-20 & CS-11 & neg. & \scriptsize\texttt{\#val20ElevenKeywordsRefused()} & \scriptsize\texttt{7.4-VAL-20-ElevenKeywordsRefused.png} \\
VAL-21 & CS-11 & neg. & \scriptsize\texttt{\#val21LongKeywordRefused()} & \scriptsize\texttt{7.4-VAL-21-LongKeywordRefused.png} \\
VAL-22 & CS-10 & pos. & \scriptsize\texttt{\#val22ContributorAccepted()} & \scriptsize\texttt{7.4-VAL-22-ContributorAccepted.png} \\
VAL-23 & CS-10 & neg. & \scriptsize\texttt{\#val23BlankContributorRefused()} & \scriptsize\texttt{7.4-VAL-23-BlankContributorRefused.png} \\
VAL-24 & CS-10, CS-11 & pos. & \scriptsize\texttt{\#val24FullLessonAccepted()} & \scriptsize\texttt{7.4-VAL-24-FullLessonAccepted.png} \\
VAL-25 & CS-10 & neg. & \scriptsize\texttt{\#val25FullLessonWithBlankTitleRefused()} & \scriptsize\texttt{7.4-VAL-25-FullLessonWithBlankTitleRefused.png} \\
\multicolumn{5}{@{}l}{\sffamily\bfseries\small LessonsCrudTestingAutomation.html} \\
CRUD-01 & CS-1 & pos. & \scriptsize\texttt{\#crud01CreateValidLesson()} & \scriptsize\texttt{7.4-CRUD-01-CreateValidLesson.png} \\
CRUD-02 & CS-1 & neg. & \scriptsize\texttt{\#crud02CreateBlankTitleRefused()} & \scriptsize\texttt{7.4-CRUD-02-CreateBlankTitleRefused.png} \\
CRUD-03 & CS-3 & pos. & \scriptsize\texttt{\#crud03ReadById()} & \scriptsize\texttt{7.4-CRUD-03-ReadById.png} \\
CRUD-04 & CS-3 & neg. & \scriptsize\texttt{\#crud04ReadUnknownId()} & \scriptsize\texttt{7.4-CRUD-04-ReadUnknownId.png} \\
CRUD-05 & CS-7 & pos. & \scriptsize\texttt{\#crud05UpdatePersists()} & \scriptsize\texttt{7.4-CRUD-05-UpdatePersists.png} \\
CRUD-06 & CS-7 & neg. & \scriptsize\texttt{\#crud06UpdateBlankTitleRefused()} & \scriptsize\texttt{7.4-CRUD-06-UpdateBlankTitleRefused.png} \\
CRUD-07 & CS-8 & pos. & \scriptsize\texttt{\#crud07Delete()} & \scriptsize\texttt{7.4-CRUD-07-Delete.png} \\
CRUD-08 & CS-8 & neg. & \scriptsize\texttt{\#crud08DeleteTwice()} & \scriptsize\texttt{7.4-CRUD-08-DeleteTwice.png} \\
CRUD-09 & CS-9 & pos. & \scriptsize\texttt{\#crud09OwnerRecognised()} & \scriptsize\texttt{7.4-CRUD-09-OwnerRecognised.png} \\
CRUD-10 & CS-9 & neg. & \scriptsize\texttt{\#crud10NonOwnerNotRecognised()} & \scriptsize\texttt{7.4-CRUD-10-NonOwnerNotRecognised.png} \\
CRUD-11 & CS-12 & pos. & \scriptsize\texttt{\#crud11And12PersistAcrossReopen()} & \scriptsize\texttt{7.4-CRUD-11-CRUD-12-PersistAcrossReopen.png} \\
CRUD-12 & CS-12 & neg. & \scriptsize\texttt{\#crud11And12PersistAcrossReopen()} & \scriptsize\texttt{7.4-CRUD-11-CRUD-12-PersistAcrossReopen.png} \\
CRUD-13 & CS-12 & pos. & \scriptsize\texttt{\#crud13And14OpenHw2Layout()} & \scriptsize\texttt{7.4-CRUD-13-CRUD-14-OpenHw2Layout.png} \\
CRUD-14 & CS-12 & neg. & \scriptsize\texttt{\#crud13And14OpenHw2Layout()} & \scriptsize\texttt{7.4-CRUD-13-CRUD-14-OpenHw2Layout.png} \\
\multicolumn{5}{@{}l}{\sffamily\bfseries\small LessonsListTestingAutomation.html} \\
LIST-01 & CS-2 & pos. & \scriptsize\texttt{\#list01ReadAllReturnsEveryLesson()} & \scriptsize\texttt{7.4-LIST-01-ReadAllReturnsEveryLesson.png} \\
LIST-02 & CS-2 & neg. & \scriptsize\texttt{\#list02EmptyDatabaseGivesEmptyList()} & \scriptsize\texttt{7.4-LIST-02-EmptyDatabaseGivesEmptyList.png} \\
LIST-03 & CS-4 & pos. & \scriptsize\texttt{\#list03SearchFindsKeyword()} & \scriptsize\texttt{7.4-LIST-03-SearchFindsKeyword.png} \\
LIST-04 & CS-4 & neg. & \scriptsize\texttt{\#list04SearchForAbsentTerm()} & \scriptsize\texttt{7.4-LIST-04-SearchForAbsentTerm.png} \\
LIST-05 & CS-5 & pos. & \scriptsize\texttt{\#list05FilterByCategory()} & \scriptsize\texttt{7.4-LIST-05-FilterByCategory.png} \\
LIST-06 & CS-5 & neg. & \scriptsize\texttt{\#list06FilterByAbsentCategory()} & \scriptsize\texttt{7.4-LIST-06-FilterByAbsentCategory.png} \\
LIST-07 & CS-6 & pos. & \scriptsize\texttt{\#list07MyLessons()} & \scriptsize\texttt{7.4-LIST-07-MyLessons.png} \\
LIST-08 & CS-6 & neg. & \scriptsize\texttt{\#list08LessonsOfUnknownUser()} & \scriptsize\texttt{7.4-LIST-08-LessonsOfUnknownUser.png} \\
LIST-09 & CS-5 & pos. & \scriptsize\texttt{\#list09CategoriesOffered()} & \scriptsize\texttt{7.4-LIST-09-CategoriesOffered.png} \\
LIST-10 & CS-4, CS-5 & pos. & \scriptsize\texttt{\#list10SearchThenFilter()} & \scriptsize\texttt{7.4-LIST-10-SearchThenFilter.png} \\
LIST-11 & CS-2 & pos. & \scriptsize\texttt{\#list11NewestFirst()} & \scriptsize\texttt{7.4-LIST-11-NewestFirst.png} \\
LIST-12 & CS-4 & pos. & \scriptsize\texttt{\#list12SearchIgnoresCase()} & \scriptsize\texttt{7.4-LIST-12-SearchIgnoresCase.png} \\
LIST-13 & CS-4 & neg. & \scriptsize\texttt{\#list13EmptySearchTerm()} & \scriptsize\texttt{7.4-LIST-13-EmptySearchTerm.png} \\
\multicolumn{5}{@{}l}{\sffamily\bfseries\small LessonsModelTestingAutomation.html} \\
MOD-01 & CS-1 & pos. & \scriptsize\texttt{\#mod01SaveNewLesson()} & \scriptsize\texttt{7.4-MOD-01-SaveNewLesson.png} \\
MOD-02 & CS-1 & neg. & \scriptsize\texttt{\#mod02SaveNewLessonBlankTitle()} & \scriptsize\texttt{7.4-MOD-02-SaveNewLessonBlankTitle.png} \\
MOD-03 & CS-2 & pos. & \scriptsize\texttt{\#mod03NoFiltersShowsAll()} & \scriptsize\texttt{7.4-MOD-03-NoFiltersShowsAll.png} \\
MOD-04 & CS-6 & neg. & \scriptsize\texttt{\#mod04MessageWhenUserHasNoLessons()} & \scriptsize\texttt{7.4-MOD-04-MessageWhenUserHasNoLessons.png} \\
MOD-05 & CS-4 & pos. & \scriptsize\texttt{\#mod05Search()} & \scriptsize\texttt{7.4-MOD-05-Search.png} \\
MOD-06 & CS-4 & neg. & \scriptsize\texttt{\#mod06SearchForAbsentTerm()} & \scriptsize\texttt{7.4-MOD-06-SearchForAbsentTerm.png} \\
MOD-07 & CS-5 & pos. & \scriptsize\texttt{\#mod07CategoryFilter()} & \scriptsize\texttt{7.4-MOD-07-CategoryFilter.png} \\
MOD-08 & CS-5 & neg. & \scriptsize\texttt{\#mod08AbsentCategory()} & \scriptsize\texttt{7.4-MOD-08-AbsentCategory.png} \\
MOD-09 & CS-6 & pos. & \scriptsize\texttt{\#mod09OnlyMine()} & \scriptsize\texttt{7.4-MOD-09-OnlyMine.png} \\
MOD-10 & CS-6 & neg. & \scriptsize\texttt{\#mod10OnlyMineForUserWithNone()} & \scriptsize\texttt{7.4-MOD-10-OnlyMineForUserWithNone.png} \\
MOD-11 & CS-9 & pos. & \scriptsize\texttt{\#mod11OwnerMayChange()} & \scriptsize\texttt{7.4-MOD-11-OwnerMayChange.png} \\
MOD-12 & CS-9 & neg. & \scriptsize\texttt{\#mod12NonOwnerMayNotChange()} & \scriptsize\texttt{7.4-MOD-12-NonOwnerMayNotChange.png} \\
MOD-13 & CS-7 & pos. & \scriptsize\texttt{\#mod13OwnerUpdate()} & \scriptsize\texttt{7.4-MOD-13-OwnerUpdate.png} \\
MOD-14 & CS-7 & neg. & \scriptsize\texttt{\#mod14UpdateBlankDescription()} & \scriptsize\texttt{7.4-MOD-14-UpdateBlankDescription.png} \\
MOD-15 & CS-8 & pos. & \scriptsize\texttt{\#mod15OwnerDelete()} & \scriptsize\texttt{7.4-MOD-15-OwnerDelete.png} \\
MOD-16 & CS-8, CS-9 & neg. & \scriptsize\texttt{\#mod16NonOwnerDelete()} & \scriptsize\texttt{7.4-MOD-16-NonOwnerDelete.png} \\
MOD-17 & CS-4 & neg. & \scriptsize\texttt{\#mod17MessageForNoSearchMatch()} & \scriptsize\texttt{7.4-MOD-17-MessageForNoSearchMatch.png} \\
MOD-18 & CS-5 & neg. & \scriptsize\texttt{\#mod18MessageForNoCategoryMatch()} & \scriptsize\texttt{7.4-MOD-18-MessageForNoCategoryMatch.png} \\
MOD-19 & CS-2 & neg. & \scriptsize\texttt{\#mod19MessageWhenNoLessonsStored()} & \scriptsize\texttt{7.4-MOD-19-MessageWhenNoLessonsStored.png} \\
MOD-20 & CS-9 & neg. & \scriptsize\texttt{\#mod20NonOwnerEdit()} & \scriptsize\texttt{7.4-MOD-20-NonOwnerEdit.png} \\
MOD-21 & CS-11 & neg. & \scriptsize\texttt{\#mod21LeadingEmptyKeywordRefused()} & \scriptsize\texttt{7.4-MOD-21-LeadingEmptyKeywordRefused.png} \\
MOD-22 & CS-11 & neg. & \scriptsize\texttt{\#mod22TrailingEmptyKeywordRefused()} & \scriptsize\texttt{7.4-MOD-22-TrailingEmptyKeywordRefused.png} \\
MOD-23 & CS-5 & neg. & \scriptsize\texttt{\#mod23ReservedCategoryRefused()} & \scriptsize\texttt{7.4-MOD-23-ReservedCategoryRefused.png} \\
MOD-24 & CS-5 & pos. & \scriptsize\texttt{\#mod24SimilarCategoryNamesAccepted()} & \scriptsize\texttt{7.4-MOD-24-SimilarCategoryNamesAccepted.png} \\
\multicolumn{5}{@{}l}{\sffamily\bfseries\small LessonsManualTestProcedures.html} \\
MAN-01 & CS-2 & neg. & \scriptsize\texttt{\#man01EmptyListMessage()} & \scriptsize\texttt{7.4-MAN-01-EmptyListMessage.png} \\
MAN-02 & CS-1 & pos. & \scriptsize\texttt{\#man02CreateSetsContributor()} & \scriptsize\texttt{7.4-MAN-02-CreateSetsContributor.png} \\
MAN-03 & CS-10 & neg. & \scriptsize\texttt{\#man03RefusedFieldIsNamed()} & \scriptsize\texttt{7.4-MAN-03-RefusedFieldIsNamed.png} \\
MAN-04 & CS-8 & pos. + neg. & \scriptsize\texttt{\#man04DeleteAsksFirst()} & \scriptsize\texttt{7.4-MAN-04-DeleteAsksFirst.png} \\
MAN-05 & CS-9 & neg. & \scriptsize\texttt{\#man05OtherUserCannotChange()} & \scriptsize\texttt{7.4-MAN-05-OtherUserCannotChange.png} \\
MAN-06 & CS-8 & pos. & \scriptsize\texttt{\#man06DeleteSeenByOtherUser()} & \scriptsize\texttt{7.4-MAN-06-DeleteSeenByOtherUser.png} \\
MAN-07 & CS-12 & pos. & \scriptsize\texttt{\#man07RestartKeepsLessons()} & \scriptsize\texttt{7.4-MAN-07-RestartKeepsLessons.png} \\
\bottomrule\end{longtable}
\end{small}

\subsection*{Other screenshots in ContributorTests\_Javadocs}
\begin{tabularx}{\linewidth}{@{}l X@{}}
\toprule
\textbf{Screenshot} & \textbf{Shows} \\
\midrule
\texttt{7.4-<suite>-class-javadoc.png} (five) & Each suite's class Javadoc: requirements covered, test data, how to read the output \\
\texttt{7.4-TP2TestRunner-class-javadoc.png}, \texttt{-main.png} & What each suite covers; how the totals are added \\
\texttt{7.4-check-helper-and-refusedWith.png} & How a Pass or Failure line is produced; why a negative test checks the refusal message \\
\texttt{7.4-TP2TestRunner-output-1} to \texttt{-5} & The console output of a real run, one image per suite and one for the totals \\
\bottomrule
\end{tabularx}

\section*{Latest run}
October 26, 2026, on the submitted source:

\begin{tabularx}{\linewidth}{@{}X r r r@{}}
\toprule
\textbf{Suite} & \textbf{Tests} & \textbf{Passed} & \textbf{Failed} \\
\midrule
LessonValidationTestingAutomation (VAL) & 25 & 25 & 0 \\
LessonsCrudTestingAutomation (CRUD) & 14 & 14 & 0 \\
LessonsListTestingAutomation (LIST) & 13 & 13 & 0 \\
LessonsModelTestingAutomation (MOD) & 24 & 24 & 0 \\
\midrule
\textbf{Total, automated} & \textbf{76} & \textbf{76} & \textbf{0} \\
\bottomrule
\end{tabularx}

To check that the tests can fail, three defects were put into a copy of the code: the HW2-table
fix removed, keywords dropped from the search, and the list order reversed. The run then reported
six failures (CRUD-13, CRUD-14, LIST-03, LIST-11, LIST-12, MOD-05) and no false passes.

The manual procedures are carried out during the TestingScreencast; their results are recorded
there.

\section*{Iteration notes (Task 6)}
\begin{tabularx}{\linewidth}{@{}l X@{}}
\toprule
\textbf{Date} & \textbf{What the tests found, and what changed} \\
\midrule
Oct 14 & VAL-12 (a 300-character description) failed on its first run. \texttt{InputLimits.checkLength}
  stops at 256 characters, so every description from 257 to 2000 was refused. Added
  \texttt{LessonInputPolicy.checkLongText} for description and resolution. \\
Oct 19 & First CRUD run: 10 of 12. CRUD-02 and CRUD-06 failed because \texttt{create} and \texttt{update}
  wrote the row without acting on \texttt{validate}'s result. Both now return the problem and write
  nothing. \\
Oct 22 & First GUI Model run: MOD-04 failed. \texttt{statusMessage} said nothing matched when only
  ``Only my lessons'' was ticked; the branch for that case now comes first. Full run: 63 of 63. \\
Oct 25 & Review before submission. (1) Saving over the \texttt{lessonsLearned} table left by an HW2
  version (columns author, body, tags) failed with ``Column DESCRIPTION not found'', and the list
  showed nothing. \texttt{LessonsLearnedDatabase} now renames such a table to
  \texttt{lessonsLearnedHW2Legacy}, keeping its rows; CRUD-13 and CRUD-14 added. (2) CRUD-11 and
  CRUD-12 only counted rows in an in-memory database, which cannot show persistence; they now close
  and reopen a database file. (3) Each test became its own documented method; negative validation
  tests now check the refusal message. (4) Added LIST-11 to LIST-13 (order, case, empty term),
  MOD-17 to MOD-20 (no-match messages, editing someone else's lesson), and MAN-01 to MAN-07.
  Result: 72 of 72. \\
Oct 26 & Moksha's review of the editor (pull request \#1) found two more defects. An empty first
  keyword (``, java'') was silently removed by \texttt{tidyKeywords} and the lesson accepted, and
  \texttt{validateKeywords} dropped a trailing empty entry (``java,''). A lesson could also be
  stored under ``All categories'', the filter's own label, and then never be shown on its own.
  Fixed in \texttt{ModelLessonEditor} and \texttt{LessonInputPolicy}, with the label now a single
  constant shared by the filter and the rule; MOD-21 to MOD-24 added. Result: 76 of 76. \\
\bottomrule
\end{tabularx}

\end{document}
